\section{Lean formalization and the ProveIt boundary}\label{sec:lean}
\subsection{Checked statements in the archive}
Five standalone source files were checked with Lean 4.32.0. Their numerical inputs and the preceding encodings allow comparison with the arithmetic statements. The axiom reports below are Lean metatheoretic dependencies. In particular they describe the proofs of the Lean statements, not extra object axioms of iPRA.

\begin{center}\small\renewcommand{\arraystretch}{1.17}
\begin{tabular}{@{}p{.21\linewidth}p{.495\linewidth}>{\raggedright\arraybackslash}p{.22\linewidth}@{}}\toprule
File & Kernel-checked content & Axiom audit of named endpoints\\\midrule
\texttt{Bounded.lean} & Least numerical witness with sentinel, its range, success, minimality and a constructive existential equivalence. Executable Hall condition and zero-sum selector. & \texttt{propext}, \texttt{Quot.sound} for search\\
\texttt{WordCount.lean} & Derivation tree of $W_n=(a+b)^n$, then a soundness proof interpreting all rules over naturals. & None for certificate. The interpreted equality uses \texttt{propext}\\
\texttt{Handshaking.lean} & Finite double-sum interchange, weighted symmetric degree sum, $\sum_v d(v)=2e$ for a graph mask, and an explicit half of the odd-degree count. & None for all four reported theorems\\
\texttt{Kirkman.lean} & Seven days with five valid triples each, one attendance per pupil per day, and one meeting for every $x<y<15$. & None for all three theorems\\
\texttt{ThueMorse.lean} & Digit recurrences, scalar PR prefix coding, odd/even overlap descent, a bounded mismatch, aperiodicity, cube-freeness and prefix balance. & \texttt{propext}, \texttt{Quot.sound} for the reported endpoints\\\bottomrule
\end{tabular}
\end{center}
The word-count certificate uses a small equational language: the base and step constructors for powers and word counts, multiplication congruence, symmetry and transitivity, distributivity, and induction on equality of two terms. Its induction premises substitute zero and successor for the sole induction variable, leaving two free letter-count parameters. Reading these constructors as the corresponding iPRA rules is a rule-by-rule interpretation; distributivity is a derived semiring law. The other standalone files prove ordinary Lean propositions. The integral/series arguments, uniform hypergraph Ramsey theorem, hook formulas, polytope results and other results outside these endpoints have the human proofs and formalization plans in their respective sections.

\subsection{Constructing derivations in the repository's ordinary proof calculus}
The inspected repository revision is \texttt{b4de6f7a61a3fcf9740373b93089830f79e342f6}. Its \texttt{ConciseIPRA.PR} represents terms as \texttt{PTm n}, formulas by \texttt{Fm PTm n}, and derivations by a judgment of the form \(\Gamma\vdash[\mathrm{PRAAx}]\varphi\). Its \texttt{IPRADerivations} library supplies arithmetic and inference steps; its research packets have separate verification statuses. The combinatorics packet directory at this revision lists Bell numbers, Catalan numbers, Fibonacci--Lucas numbers, and finite Ramsey theory. For parameterized bijections use the following order of work.

\begin{enumerate}
\item \textbf{Define the representation terms.} Introduce powers, quotient, parity, masked size, lookup, padding and set difference. Express each table update with composition or primitive recursion. Use the pairing projections to pass multiple parameters to a bounded sum. For a new object give a Boolean validity term, domain bound and canonical decoding. Prove the two radix laws in Section~\ref{sec:found} using division with remainder; then iteration supplies lookup-after-write and prefix reconstruction. Translate the hypotheses on graphs, families, moves and covers into bounded Boolean terms by Lemma~\ref{lem:search}.
\item \textbf{Prove counting once.} Use \texttt{sum\_zero\_elim} to extract pointwise validity from zero violation totals; \texttt{Search.search\_found} to extract a successful witness; \texttt{strong\_ind} when all earlier indices are needed. The repository theorem \texttt{FiniteBijection.finite\_bijection} asks for four object derivations: the forward and reverse maps preserve bounds and valid objects, and each round trip is the identity on its own valid domain. It then identifies the two indicator sums. The prefix arguments in our bijective proofs establish precisely these obligations. When two valid encodings can denote the same object, restrict to canonical representatives or prove invariance and count the explicitly sized fibers.
\item \textbf{Illustration with Pr\"ufer.} Use one common numerical parameter packing $n$, initial graph and current live mask as needed. The source predicate tests a simple tree on $[n]$; its bound is $2^{\binom n2}$. The target predicate tests a length-$n-2$ base-$\max(2,n)$ word and $n\ge2$, bounded by $\max(2,n)^{n-2}$. Each encoder step chooses the first vertex with live degree one, emits its neighbor and erases it; the decoder takes the first live label absent from the unprocessed suffix. In the decoder's remaining tree the degree of each live label is one plus its number of suffix occurrences. The final two live labels have degree one. Adding the deleted edge proves the degree formula backwards and recovers the exact leaf choice forwards. Use the four resulting judgments as hypotheses to the bijection theorem; then replace the count of words by its radix power. Section~\ref{sec:graphs} proves the invariants independently of the representation.
\item \textbf{Illustration with variable-rank Ramsey.} For fixed final parameters $(h,c,m)$ first form $N=H_h(c,m)$. Write $B(j)$ for the sum, over rank-$j$ color tables on $H_j$ points, of indicators that all $m$-masks fail. Encode invalid table entries as zero contributions. The induction matrix can be $B(j)=0$ guarded by $1\le j\le h$ with $h,c,m,N$ as parameters. The induction step builds the pivot list, records the thinned colors and forms one smaller coloring. Apply the sum-zero rule to that coloring, recover its least homogeneous mask and lift it to the pivot labels. The rank proof in Section~\ref{sec:graphs} bounds every table in this use. The repository's \texttt{RamseyBounds.neighbor\_split} and \texttt{pascal\_bound} cover arithmetic steps for the graph bound; they have narrower conclusions than this new uniform hypergraph statement.
\item \textbf{Close each theorem.} Replace subtraction of signed quantities by addition of the positive and negative registers on opposite sides. Use only positive denominators when clearing a ratio. Prove the Boolean implication $\mathrm{Valid}\Rightarrow\mathrm{Check}$ for free parameter terms by these bounded inductions, turn its truth equation into the object implication, and universally close the parameters. Where the mathematical statement has an existential conclusion, introduce the numerical selector with its success proof. For analytic estimates the precision index is another numeric parameter and the witnesses are rational endpoints or mesh sizes. The final target is a closed object derivation \( []\vdash[\mathrm{PRAAx}]\varphi\).
\end{enumerate}
For exact reusable names, \texttt{lean/ProveItEndpoints.lean} lists references to the pinned modules. For example \texttt{Catalan.cat\_spec} supplies $(n+1)C_n=\binom{2n}{n}$ and \texttt{Vandermonde.vandermonde} uses a guarded binomial coefficient that vanishes above its row. The names and their source signatures were checked against the repository. The optional import file has not been built in the full ProveIt workspace during this study. The five standalone files were compiled separately; the transcript accompanies the archive.

\subsection{Source scope and further questions}
The starting article was read at revision 1384635198. The file \texttt{sources/direct-links.tsv} records its English article links from the introduction through ``See also,'' with their anchor text and enclosing section. \texttt{sources/theorem-map.tsv} matches result groups to these entry points. For example, the triangle-free maximum and the questions about intersecting and inclusion-incomparable families occur in the overview; the twelvefold enumeration and the two tableau counts are developed at immediate destinations. The group-action, zero-sum and Markov-tree arguments also supply consequences along these topics. The coverage is selective. Infinite Ramsey for arbitrary set-valued colorings, the full four-color theorem, and an infinite-volume two-dimensional Ising solution lie outside the proofs given here. The finite Ising identities in Section~\ref{sec:models} specify their graph and boundary conditions.

Three next questions follow directly from the proofs. Can the entire mask-level Pr\"ufer bijection be exported to \texttt{Prf PRAAx} with a single reusable deletion lemma? Can the two-cut sliding-window argument be replaced by an explicit iPRA parity certificate for arbitrary numbers of bead colors and the optimal cut bound? How much of the uniform Ramsey proof can share a single object-language lemma for lifting a homogeneous set through an order-preserving embedding?

\begin{thebibliography}{9}\small
\bibitem{overview} Wikipedia contributors. \emph{Combinatorics}, revision 1384635198. \href{https://en.wikipedia.org/w/index.php?title=Combinatorics\&oldid=1384635198}{Article snapshot}.
\bibitem{ipra} V. Reshetnikov and contributors. \emph{ProveIt}, Logic/PrimitiveRecursiveArithmetic. \href{https://github.com/VladimirReshetnikov/ProveIt/tree/b4de6f7a61a3fcf9740373b93089830f79e342f6/Logic/PrimitiveRecursiveArithmetic}{Pinned source tree} and \href{https://github.com/VladimirReshetnikov/ProveIt/tree/b4de6f7a61a3fcf9740373b93089830f79e342f6/Logic/PrimitiveRecursiveArithmetic/Derivations/Lean/IPRADerivations/Combinatorics}{individual derivations}.
\bibitem{alon} N. Alon. Combinatorial Nullstellensatz. \emph{Combinatorics, Probability and Computing} 8 (1999), 7--29. \href{https://www.math.tau.ac.il/~nogaa/PDFS/null2.pdf}{Author's paper}.
\bibitem{sokal} A. D. Sokal. The multivariate Tutte polynomial (alias Potts model) for graphs and matroids. \emph{Surveys in Combinatorics 2005}, 173--226. \href{https://arxiv.org/abs/math/0503607}{Author's preprint}.
\bibitem{thue} J.-P. Allouche and J. Shallit. The ubiquitous Prouhet--Thue--Morse sequence. In \emph{Sequences and Their Applications}, 1999, 1--16. \href{https://cs.uwaterloo.ca/~shallit/Papers/ubiq15.pdf}{Authors' paper}.
\end{thebibliography}
